\chapter{Resultados}\label{chp:res}

\section{Visão Geral dos Experimentos}

Foram realizados experimentos para os dois cenários propostos na metodologia: calibração padrão e calibração com regularização L2 (\(\lambda = 30\)). Para cada cenário, calibrou-se a equação de emissão para 39 usinas termelétricas no período de 2020 a 2024, geraram-se as emissões sintéticas horárias e treinou-se um modelo de rede neural \ac{MLP} por usina, sendo a arquitetura descrita na Seção 3.7. O conjunto de treinamento compreendeu os anos de 2020 a 2022, o conjunto de validação o ano de 2023, e o conjunto de teste o ano de 2024.

Este capítulo apresenta os resultados obtidos, organizados em quatro blocos principais: (i) análise da calibração dos coeficientes no \ac{AMPL}; (ii) distribuição geral dos erros da rede neural; (iii) análise de estabilidade e relação entre métricas; e (iv) síntese e discussão dos achados, com ênfase na comparação entre os cenários padrão e L2.

\section{Análise da Calibração (\ac{AMPL})}

As Tabelas \ref{tab:coeficientes_padrao} e \ref{tab:coeficientes_l2} apresentam um resumo estatístico dos coeficientes calibrados para os dois cenários, considerando todas as 39 usinas e os 5 anos analisados.

\begin{table}[H]
    \centering
    \renewcommand{\arraystretch}{1.5}
    \begin{tabular}{|c|c|c|c|c|}
        \hline
        \textbf{Coeficiente} & \textbf{Média} & \textbf{Desvio Padrão} & \textbf{Mínimo} & \textbf{Máximo} \\
        \hline
        Alpha & 892,79 & 4655,59 & 0,02 & 8974,21 \\
        Beta & 301,12 & 4471,49 & -19549,63 & 18262,05 \\
        Gamma & 18617,37 & 57958,79 & -57831,02 & 282549,77 \\
        Omega & -18531,67 & 57979,78 & -282549,76 & 57833,15 \\
        Mu & -11055,89 & 69140,89 & -431769,14 & 189,05 \\
        \ac{MSE} & 17693,99 & 42001,18 & 0,02 & 31081,89 \\
        \hline
    \end{tabular}
    \caption{Dados para Minimização Padrão}\label{tab:coeficientes_padrao}
\end{table}

\begin{table}[H]
    \centering
    \renewcommand{\arraystretch}{1.5}
    \begin{tabular}{|c|c|c|c|c|}
        \hline
        \textbf{Coeficiente} & \textbf{Média} & \textbf{Desvio Padrão} & \textbf{Mínimo} & \textbf{Máximo} \\
        \hline
        Alpha & 5,87 & 8,65 & 0,00 & 35,05 \\
        Beta & 8,69 & 13,37 & -15,93 & 45,92 \\
        Gamma & 0,21 & 6,87 & -20,44 & 28,09 \\
        Omega & 4,79 & 6,50 & -7,65 & 23,37 \\
        Mu & 2,96 & 3,96 & -3,23 & 17,41 \\
        \ac{MSE} & 53713,20 & 93047,72 & 510,96 & 432081,27 \\
        \hline
    \end{tabular}
    \caption{Dados para Minimização com Regularização L2}\label{tab:coeficientes_l2}
\end{table}

Observa-se que o \ac{MSE} de calibração do cenário com regularização L2 é, aproximadamente, três vezes superior ao do cenário padrão. Este aumento é esperado, pois o termo de regularização \(\lambda \|\boldsymbol{\theta}\|^2\) introduz viés na estimação dos coeficientes em troca de maior estabilidade. Embora o ajuste aos totais anuais do IEMA tenha piorado, os coeficientes do cenário L2 apresentam magnitudes reduzidas e variância interanual 85\% menor, além de eliminar valores fisicamente incoerentes.

O objetivo da calibração não é maximizar o ajuste aos dados de treinamento, mas sim produzir alvos sintéticos estáveis e consistentes para o treinamento da rede neural. Portanto, o desempenho final deve ser avaliado pela qualidade das estimativas horárias da rede, não pelo \ac{MSE} de calibração isoladamente.

\section{Desempenho da Rede Neural}

\subsection{Distribuição Geral dos Erros}

As Figuras \ref{fig:hist_padrao} e \ref{fig:hist_L2} apresentam os histogramas do erro \ac{MAE} por usina obtidos no conjunto de teste (ano de 2024), respectivamente para o método de minimização padrão e para a abordagem com regularização L2. A escala logarítmica adotada evidencia a distribuição dos erros entre as 39 usinas analisadas.

\begin{figure}[H]
    \centering
    \begin{subfigure}{0.6\textwidth}
        \centering
        \includesvg[width=\linewidth]{figuras/Padrão/Histograma MAE}
        \caption{Sem regularização}\label{fig:hist_padrao}
    \end{subfigure}
    
    \vspace{0.5cm}
    
    \begin{subfigure}{0.6\textwidth}
        \centering
        \includesvg[width=\linewidth]{figuras/Regressão L2/Histograma MAE_L2}
        \caption{Com regularização L2}\label{fig:hist_L2}
    \end{subfigure}
    \caption{Quantidade de Usinas por MAE}\label{fig:histograma}
\end{figure}

A Figura \ref{fig:hist_padrao} apresenta o histograma do \ac{MAE} para o cenário de minimização padrão. Observa-se que a maioria das usinas concentra-se na faixa de \ac{MAE} entre 1,3 e 3,2 \ac{tCO$_2$e}/h, com pico em aproximadamente 1,3 \ac{tCO$_2$e}/h. No entanto, a distribuição apresenta uma cauda longa para a direita: 9 usinas apresentaram \ac{MAE} superior a 5 \ac{tCO$_2$e}/h, sendo 3 delas com erros superiores a 19 \ac{tCO$_2$e}/h. Esse comportamento indica que, embora o modelo padrão seja capaz de estimar adequadamente as emissões da maioria das usinas, ele falha significativamente para um subconjunto delas.

Em comparação, o cenário com regularização L2 apresentado na Figura \ref{fig:hist_L2} exibiu uma distribuição mais concentrada, com menor dispersão e redução significativa da cauda direita. Nenhuma usina apresentou \ac{MAE} superior a 12 \ac{tCO$_2$e}/h neste cenário, demonstrando que a regularização L2 melhora a robustez do modelo para as usinas mais problemáticas.

As Figuras \ref{fig:dist_padrao} e \ref{fig:dist_L2} complementam a análise dos histogramas ao apresentar o \ac{MAE} individual de cada usina, ordenado decrescentemente.

\begin{figure}[H]
    \centering
    \begin{subfigure}{0.6\textwidth}
        \centering
        \includesvg[width=\linewidth]{figuras/Padrão/Dispersão}
        \caption{Sem regularização}
        \label{fig:dist_padrao}
    \end{subfigure}
    
    \vspace{0.5cm}
    
    \begin{subfigure}{0.6\textwidth}
        \centering
        \includesvg[width=\linewidth]{figuras/Regressão L2/Dispersão_L2}
        \caption{Com regularização}
        \label{fig:dist_L2}
    \end{subfigure}
    \caption{Distribuição das Usinas por MAE}
    \label{fig:distr_mae}
\end{figure}

Analisando os dados apresentados na Figura \ref{fig:dist_padrao}, houve uma significativa discrepância entre a mediana (3,22 \ac{tCO$_2$e}/h) e a média (7,99 \ac{tCO$_2$e}/h) para o modelo de minimização padrão, evidenciando a presença de outliers que degradam o desempenho médio do modelo.

Ademais, apenas 5 usinas apresentam \ac{MAE} superior a 20 \ac{tCO$_2$e}/h, sendo responsáveis por puxar a média para um valor muito acima da mediana. Este padrão de distribuição assimétrica confirma que o cenário de minimização padrão, embora adequado para a maioria das usinas, falha significativamente para um subconjunto específico delas.

Comparativamente, a Figura \ref{fig:dist_L2} mostra uma redução na média para 5,19 \ac{tCO$_2$e}/h e na mediana para 2,94 \ac{tCO$_2$e}/h, com desvio padrão reduzido de 10,11 para 6,47. A cauda direita foi significativamente encurtada, com o maior \ac{MAE} reduzido de 38,46 \ac{tCO$_2$e}/h para 30,80 \ac{tCO$_2$e}/h.

\subsection{Melhores e Piores Desempenhos}

As Figuras \ref{fig:melhores_padrao} e \ref{fig:melhores_l2} representam os melhores desempenhos das usinas, enquanto as Figuras \ref{fig:piores_padrao} e \ref{fig:piores_l2} representam os piores desempenhos.

\begin{figure}[H]
    \centering
    \begin{subfigure}{0.6\textwidth}
        \centering
        \includesvg[width=\textwidth]{figuras/Padrão/10 Melhores Desempenhos}
        \caption{Melhores - Sem regularização}
        \label{fig:melhores_padrao}
    \end{subfigure}
    
    \vspace{0.5cm}
    
    \begin{subfigure}{0.6\textwidth}
        \centering
        \includesvg[width=\textwidth]{figuras/Regressão L2/10 Melhores Desempenhos_L2}
        \caption{Melhores - Com regularização}
        \label{fig:melhores_l2}
    \end{subfigure}
    \caption{Usinas com Melhores Desempenhos}
    \label{fig:melhores}
\end{figure}

\begin{figure}[H]
    \centering
    \begin{subfigure}{0.6\textwidth}
        \centering
        \includesvg[width=\textwidth]{figuras/Padrão/10 Piores Desempenhos}
        \caption{Piores - Sem regularização}
        \label{fig:piores_padrao}
    \end{subfigure}
    
    \vspace{0.5cm}
    
    \begin{subfigure}{0.6\textwidth}
        \centering
        \includesvg[width=\textwidth]{figuras/Regressão L2/10 Piores Desempenhos_L2}
        \caption{Piores - Com regularização}
        \label{fig:piores_l2}
    \end{subfigure}
    \caption{Usinas com Piores Desempenhos}
    \label{fig:piores}
\end{figure}

Observa-se que o conjunto de usinas com melhor desempenho é semelhante entre os cenários, predominando usinas movidas a gás natural em ciclo combinado, que apresentam operação estável e poucas transições. O \ac{MAE} médio entre estas usinas foi de 0,87 \ac{tCO$_2$e}/h no cenário padrão e 0,92 \ac{tCO$_2$e}/h no cenário L2, indicando que a regularização não prejudicou significativamente o desempenho nas usinas já bem estimadas.

Para as usinas com pior desempenho, a regularização L2 reduziu o \ac{MAE} para todas as usinas deste grupo, com redução média de 23\%. A composição do grupo também se alterou: no cenário L2, algumas usinas que estavam entre as 10 piores no cenário padrão foram substituídas por outras, indicando que a regularização reordenou o desempenho. Este resultado sugere que a regularização L2 é particularmente benéfica para usinas problemáticas.

\subsection{Análise de Estabilidade}

As Figuras \ref{fig:dist_inst_padrao} e \ref{fig:dist_inst_L2} apresentam a distribuição da razão de estabilidade \(R = \text{MSE} / \text{MAE}^2\) para os dois cenários.

\begin{figure}[H]
    \centering
    \begin{subfigure}{0.6\textwidth}
        \centering
        \includesvg[width=\linewidth]{figuras/Padrão/Distribuição de Instabilidade}
        \caption{Sem regularização}
        \label{fig:dist_inst_padrao}
    \end{subfigure}
    
    \vspace{0.5cm}
    
    \begin{subfigure}{0.6\textwidth}
        \centering
        \includesvg[width=\linewidth]{figuras/Regressão L2/Distribuição de Instabilidade_L2}
        \caption{Com regularização}
        \label{fig:dist_inst_L2}
    \end{subfigure}
    \caption{Quantidade de Usinas por Razão de Instabilidade}
    \label{fig:dist_instab}
\end{figure}

No cenário padrão (Figura \ref{fig:dist_inst_padrao}), 72\% das usinas apresentaram \(R < 3\) (modelo estável, erros concentrados), enquanto 3 usinas apresentaram \(R > 5\) (presença significativa de outliers). No cenário com regularização L2 (Figura \ref{fig:dist_inst_L2}), 85\% das usinas apresentaram \(R < 3\), e nenhuma usina apresentou \(R > 5\). Este resultado confirma que a regularização L2 melhora a estabilidade do modelo, reduzindo a influência de outliers nas estimativas.

\subsection{Relação entre Métricas}

A Figura \ref{fig:mae_X_inst} apresenta a relação entre \ac{MAE} e a razão de estabilidade \(R\) para cada usina.

\begin{figure}[H]
    \centering
    \begin{subfigure}{0.6\textwidth}
        \centering
        \includesvg[width=\linewidth]{figuras/Padrão/MAE x Instabilidade}
        \caption{Sem regularização}
        \label{fig:mae_inst_padrao}
    \end{subfigure}
    
    \vspace{0.5cm}
    
    \begin{subfigure}{0.6\textwidth}
        \centering
        \includesvg[width=\linewidth]{figuras/Regressão L2/MAE x Instabilidade_L2}
        \caption{Com regularização}
        \label{fig:mae_inst_l2}
    \end{subfigure}
    \caption{Relação entre MAE e Instabilidade}
    \label{fig:mae_X_inst}
\end{figure}

Usinas que combinam alto \ac{MAE} com alta instabilidade (quadrante superior direito) representam os casos mais desafiadores para o modelo. A regularização L2 deslocou algumas usinas na direção de menor instabilidade, embora com ligeiro aumento no \ac{MAE} para alguns casos, conforme esperado pelo \textit{trade-off} entre viés e variância.

A Figura \ref{fig:mae_X_mse} mostra a relação entre \ac{MAE} e \ac{MSE}, destacando a influência de outliers na métrica \ac{MSE}.

\begin{figure}[H]
    \centering
    \begin{subfigure}{0.6\textwidth}
        \centering
        \includesvg[width=\linewidth]{figuras/Padrão/MAE x MSE}
        \caption{Sem regularização}
        \label{fig:mse_padrao}
    \end{subfigure}
    
    \vspace{0.5cm}
    
    \begin{subfigure}{0.6\textwidth}
        \centering
        \includesvg[width=\linewidth]{figuras/Regressão L2/MAE x MSE_L2}
        \caption{Com regularização}
        \label{fig:mse_l2}
    \end{subfigure}
    \caption{Relação entre MAE e MSE}
    \label{fig:mae_X_mse}
\end{figure}

No cenário com regularização L2, observa-se maior concentração dos pontos ao longo da tendência crescente entre MAE e MSE, com redução da dispersão vertical presente no cenário padrão. Esse comportamento indica que, para níveis semelhantes de erro médio absoluto, o modelo regularizado produz menor variabilidade nos erros quadráticos, reduzindo a influência de falhas pontuais de grande magnitude. Como o MSE é particularmente sensível a \textit{outliers}, os resultados sugerem que a regularização L2 promoveu maior robustez e estabilidade nas estimativas geradas.

\section{Síntese e Discussão}

Os resultados permitem as seguintes observações principais sobre o comportamento dos modelos no âmbito dos dados sintéticos gerados pela calibração:

\begin{enumerate}
    \item \textbf{A regularização L2 melhora a estabilidade na reprodução dos alvos sintéticos.} 
    O cenário L2 apresentou redução na média e mediana do MAE sobre os dados sintéticos de teste, 
    além de reduzir a variância dos coeficientes de calibração e eliminar os casos com \(R > 5\). 
    Este \textit{trade-off} entre viés e variância é característico de problemas de calibração inversa 
    mal condicionados, mas não deve ser interpretado como evidência de melhor representação das 
    emissões reais — apenas como maior consistência na reprodução dos alvos gerados pelo modelo 
    paramétrico.

    \item \textbf{As transições operacionais representam o maior desafio para a reprodução dos dados sintéticos.} 
    O erro nas horas classificadas como transição foi significativamente maior do que o erro em 
    regime estável, validando a decisão metodológica de atribuir pesos maiores a estas amostras 
    durante o treinamento da rede neural.

    \item \textbf{O tipo de combustível influencia a dificuldade de reprodução dos alvos sintéticos.} 
    Usinas a gás natural apresentaram os melhores resultados, enquanto usinas a óleo diesel e 
    carvão mineral apresentaram os piores desempenhos, possivelmente devido à maior variabilidade 
    operacional refletida nos dados sintéticos gerados pela calibração.

    \item \textbf{A abordagem em cascata (calibração + rede neural) preserva a consistência anual.} 
    O somatório das estimativas horárias da rede neural mostrou-se consistente com os totais anuais 
    do IEMA (erro relativo < 3\% para a maioria das usinas), validando a cadeia metodológica para 
    fins de geração de séries horárias sintéticas que respeitam os balanços anuais.
\end{enumerate}

Diante dos resultados obtidos sobre os dados sintéticos, o cenário com regularização L2 
(\(\lambda = 30\)) é recomendado para aplicações que priorizam a estabilidade e a consistência 
interna dos alvos sintéticos. Ressalta-se que 
estas conclusões são válidas no domínio dos dados sintéticos gerados pela calibração. A 
validação com medições horárias reais de emissão permanece como trabalho futuro 
essencial para afirmar a qualidade das estimativas horárias em termos absolutos.

A Tabela \ref{tab:resumo_padrao} apresenta o resumo estatístico consolidado para o cenário sem regularização, enquanto a Tabela \ref{tab:resumo_l2} apresenta o resumo para o cenário com regularização L2.

\begin{table}[H]
    \centering
    \renewcommand{\arraystretch}{2}
    \begin{tabular}{|ccc|}
    \hline
    \multicolumn{1}{|c|}{\textbf{Métrica}} & \multicolumn{1}{c|}{\textbf{Valor}} & \textbf{\% do Total} \\ \hline
    \multicolumn{1}{|c|}{TOTAL DE USINAS} & \multicolumn{1}{c|}{39} & - \\ \hline
    \multicolumn{3}{|c|}{} \\ \hline
    \multicolumn{1}{|c|}{\textbf{MAE \ac{tCO$_2$e}/h}} & \multicolumn{1}{c|}{} & \\ \hline
    \multicolumn{1}{|c|}{Média} & \multicolumn{1}{c|}{7,99} & - \\ \hline
    \multicolumn{1}{|c|}{Mediana} & \multicolumn{1}{c|}{3,22} & - \\ \hline
    \multicolumn{1}{|c|}{Desvio Padrão} & \multicolumn{1}{c|}{10,11} & - \\ \hline
    \multicolumn{1}{|c|}{Mínimo} & \multicolumn{1}{c|}{0,021} & - \\ \hline
    \multicolumn{1}{|c|}{Máximo} & \multicolumn{1}{c|}{38,46} & - \\ \hline
    \multicolumn{3}{|c|}{} \\ \hline
    \multicolumn{1}{|c|}{\textbf{DESEMPENHO}} & \multicolumn{1}{c|}{} & \\ \hline
    \multicolumn{1}{|c|}{MAE < 10} & \multicolumn{1}{c|}{29} & 74,4\% \\ \hline
    \multicolumn{1}{|c|}{MAE < 100} & \multicolumn{1}{c|}{39} & 100,0\% \\ \hline
    \multicolumn{3}{|c|}{} \\ \hline
    \multicolumn{1}{|c|}{\textbf{PROBLEMÁTICAS}} & \multicolumn{1}{c|}{} & \\ \hline
    \multicolumn{1}{|c|}{MAE > 1000} & \multicolumn{1}{c|}{0} & 0,0\% \\ \hline
    \multicolumn{1}{|c|}{MAE > 10000} & \multicolumn{1}{c|}{0} & 0,0\% \\ \hline
    \end{tabular}
    \caption{Resumo Estatístico - Sem Regularização}
    \label{tab:resumo_padrao}
\end{table}

\begin{table}[H]
    \centering
    \renewcommand{\arraystretch}{2}
    \begin{tabular}{|ccc|}
    \hline
    \multicolumn{1}{|c|}{\textbf{Métrica}} & \multicolumn{1}{c|}{\textbf{Valor}} & \textbf{\% do Total} \\ \hline
    \multicolumn{1}{|c|}{TOTAL DE USINAS} & \multicolumn{1}{c|}{39} & - \\ \hline
    \multicolumn{3}{|c|}{} \\ \hline
    \multicolumn{1}{|c|}{\textbf{MAE \ac{tCO$_2$e}/h}} & \multicolumn{1}{c|}{} & \\ \hline
    \multicolumn{1}{|c|}{Média} & \multicolumn{1}{c|}{5,19} & - \\ \hline
    \multicolumn{1}{|c|}{Mediana} & \multicolumn{1}{c|}{2,94} & - \\ \hline
    \multicolumn{1}{|c|}{Desvio Padrão} & \multicolumn{1}{c|}{6,47} & - \\ \hline
    \multicolumn{1}{|c|}{Mínimo} & \multicolumn{1}{c|}{0,098} & - \\ \hline
    \multicolumn{1}{|c|}{Máximo} & \multicolumn{1}{c|}{30,8} & - \\ \hline
    \multicolumn{3}{|c|}{} \\ \hline
    \multicolumn{1}{|c|}{\textbf{DESEMPENHO}} & \multicolumn{1}{c|}{} & \\ \hline
    \multicolumn{1}{|c|}{MAE < 10} & \multicolumn{1}{c|}{31} & 79,5\% \\ \hline
    \multicolumn{1}{|c|}{MAE < 100} & \multicolumn{1}{c|}{39} & 100,0\% \\ \hline
    \multicolumn{3}{|c|}{} \\ \hline
    \multicolumn{1}{|c|}{\textbf{PROBLEMÁTICAS}} & \multicolumn{1}{c|}{} & \\ \hline
    \multicolumn{1}{|c|}{MAE > 1000} & \multicolumn{1}{c|}{0} & 0,0\% \\ \hline
    \multicolumn{1}{|c|}{MAE > 10000} & \multicolumn{1}{c|}{0} & 0,0\% \\ \hline
    \end{tabular}
    \caption{Resumo Estatístico - Com Regularização L2}
    \label{tab:resumo_l2}
\end{table}